\section*{Resumen}
\addcontentsline{toc}{section}{Resumen}
Este trabajo plantea una revisión comparativa de la evaluación de calidad de imágenes faciales (\textit{Face Image Quality Assessment}, FIQA), con especial atención a los modelos de visión-lenguaje (\textit{Vision-Language Models}, VLMs). La pregunta principal es si estos modelos permiten estimar y explicar la calidad de una imagen facial de forma alineada con su utilidad para el reconocimiento biométrico, más allá de su calidad perceptual.

La revisión se organizará desde los fundamentos de calidad de imagen y los modelos multimodales hasta los métodos FIQA especializados y las propuestas basadas en VLMs. Se analizarán los protocolos y resultados publicados, atendiendo a la utilidad biométrica, las explicaciones, la sensibilidad a las instrucciones textuales y el coste computacional. Esta primera versión fija la estructura y los criterios de análisis; todavía no presenta resultados ni conclusiones de la revisión.

\textbf{Palabras clave:} calidad facial, FIQA, modelos de visión-lenguaje, reconocimiento facial, utilidad biométrica, explicabilidad.
\pendiente{Actualizar el resumen con los hallazgos de la revisión una vez completada.}
